\chapter{Conclusion and Future Work}
\label{ch:conclusion}

We have thus built a pipeline which takes a C-like sequential packet
parsing code and produces a state machine as a P4 parser using program
synthesis. The source code is available at~\cite{parsercompiler26_repo}.

\section{Limitations}
\label{sec:limitations}

Our design has a few limitations:
\begin{enumerate}
    \item The DSL is very restrictive. It only allows extracts,
    no lookaheads, no conditions involving \texttt{\&\&} or
    \texttt{||}, and no cross-packet state. All of these would be
    desirable in some cases for packet processing.

    \item We have ignored the difference between the byte order in
    network packets and the way structs are stored in memory.

    \item Bitmask handling doesn't scale well. See the Class~3 case
    in Section~\ref{sec:field-size-min} and the ``alternative we
    considered'' subsection in the IRGenerator.

    \item The state machine produced doesn't necessarily have the
    minimum number of states.
\end{enumerate}

\section{Future Work}

To build on this work:
\begin{enumerate}
    \item Extend the design to overcome the limitations above.

    \item Develop a more efficient method to handle bitmasks and automate the drop-and-stitch optimisation.
    
    \item Investigate better synthesis analyses to improve
    performance and scalability.

    \item Use these ideas for further analyses to enable vector
    processing, in particular handling the control flow divergence
    that arises due to parsing.
\end{enumerate}